\section*{Problems}

\subsection*{Problem statements}

% Remember
\begin{problema}
	\label{prob:4e01dbd}
	A continuous random variable $X$ has triangular PDF $f(x) = 2x$ for $0 \le x \le 1$ and $f(x) = 0$ otherwise.
	\begin{enumerate}
		\item Compute the expectation $\E[X] = \int_{0}^{1} x \cdot 2x\,dx$.
		\item Compute $\E[X^{2}]$ and verify the variance $\Var(X) = \E[X^{2}] - (\E[X])^{2}$.
		\item Compute the median of $X$ and compare it with $\E[X]$.
	\end{enumerate}
\end{problema}

% Understand
\begin{problema}
	\label{prob:f43c638}
	Let $X \sim \text{Exp}(\lambda = 0.5)$ with PDF $f(x) = 0.5 e^{-0.5 x}$ for $x \ge 0$.
	\begin{enumerate}
		\item Compute $\E[X] = 1/\lambda$ and $\Var(X) = 1/\lambda^{2}$ using the integral definitions.
		\item Verify the \emph{memoryless} property: $\E[X - a \mid X > a] = \E[X]$ for every $a > 0$, and explain intuitively what this property means for a waiting process modeled with this distribution.
	\end{enumerate}
\end{problema}

% Apply
\begin{problema}
	\label{prob:7b147c4}
	In a physics experiment, a quantity $X = \mu + \epsilon$ is measured, where $\mu$ is the true value and $\epsilon \sim N(0, \sigma^{2})$ is measurement error. There are $n$ independent measurements $X_{1}, \dots, X_{n}$.
	\begin{enumerate}
		\item Compute $\E[\bar{X}]$ and $\Var(\bar{X})$, where $\bar{X} = (1/n) \sum X_{i}$ is the sample mean.
		\item Prove that the standard error of the estimate is $\sigma/\sqrt{n}$ and decreases with $\sqrt{n}$.
		\item For $\sigma = 0.1$, determine $n$ such that $\sigma/\sqrt{n} < 0.01$ (error $< 10\%$ of the true value).
	\end{enumerate}
\end{problema}

% Analyze
\begin{problema}
	\label{prob:de8d740}
	Let $X$ be a continuous random variable with PDF $f(x)$. Use the law of total expectation to prove the \emph{variance decomposition formula}:
	$$\Var(X) = \E[\Var(X \mid Y)] + \Var(\E[X \mid Y])$$
	where $Y$ is another random variable. This identity is fundamental in analysis of variance (ANOVA) and hierarchical models.
\end{problema}

% Evaluate
\begin{problema}
	\label{prob:9d4a41b}
	For an exponential random variable $X \sim \text{Exp}(\lambda = 1)$:
	\begin{enumerate}
		\item Compute the skewness coefficient $\gamma_{1} = \E[(X - \mu)^{3}]/\sigma^{3}$ and the kurtosis $\gamma_{2} = \E[(X - \mu)^{4}]/\sigma^{4} - 3$.
		\item Compare these values with those of the standard Normal distribution ($\gamma_{1} = 0$, $\gamma_{2} = 0$) and interpret what they imply for the shape of the distribution (tail direction and weight of extreme values).
	\end{enumerate}
\end{problema}

% Create
\begin{problema}
	\label{prob:5c186f2}
	Design an original problem applying linearity of expectation: an investment portfolio consists of three assets with random returns $R_1$, $R_2$, $R_3$ (define their expected values and the portfolio weights $w_1, w_2, w_3$, with $w_1+w_2+w_3=1$). Compute the total expected return of the portfolio $\E[w_1R_1+w_2R_2+w_3R_3]$, and explain why this linearity property is valid even if the returns of the three assets are correlated with one another.
\end{problema}

\subsection*{Detailed solutions}

% Remember
\begin{solproblema}[prob:4e01dbd]
	\begin{enumerate}
		\item $\E[X] = \int_{0}^{1} x \cdot 2x\,dx = \int_{0}^{1} 2x^{2}\,dx = \left[ \dfrac{2x^{3}}{3} \right]_{0}^{1} = \dfrac{2}{3}$.
		\item $\E[X^{2}] = \int_{0}^{1} 2x^{3}\,dx = \dfrac{1}{2}$. Therefore, $\Var(X) = \dfrac{1}{2} - \dfrac{4}{9} = \dfrac{1}{18} \approx 0.0556$.
		\item The median $m$ satisfies $F(m) = m^{2} = 0.5$, giving $m = 1/\sqrt{2} \approx 0.7071$. Since $\E[X] = 2/3 \approx 0.6667 < m$, the mean is less than the median because of the asymmetry of the triangular distribution: the left tail contains less mass but extends comparatively, pulling the mean below the median.
	\end{enumerate}
\end{solproblema}

% Understand
\begin{solproblema}[prob:f43c638]
	\begin{enumerate}
		\item By integration by parts with $u=x$, $dv=0.5e^{-0.5x}dx$:
		\begin{align*}
			\E[X] = \int_{0}^{\infty} x \cdot 0.5 e^{-0.5x}\,dx = \left[-xe^{-0.5x}\right]_0^\infty + \int_0^\infty e^{-0.5x}\,dx = 0+2 = 2 = 1/\lambda.
		\end{align*}
		Similarly, $\E[X^2] = 8 = 1/\lambda^2$, so $\Var(X) = 8-4=4=1/\lambda^2$.
		\item Conditional on $X>a$, $X-a$ follows $\text{Exp}(\lambda)$, so $\E[X-a\mid X>a] = 1/\lambda = \E[X]$. Intuitively, this means that an exponential waiting process \emph{does not age}: if one has already waited $a$ units of time without the event occurring, the expected additional waiting time is exactly the same as at the beginning, as if the clock had restarted. This is why the exponential models processes without fatigue or accumulated wear (for example, interarrival times in queueing theory).
	\end{enumerate}
\end{solproblema}

% Apply
\begin{solproblema}[prob:7b147c4]
	\begin{enumerate}
		\item Each measurement is $X_i=\mu+\epsilon_i$ with independent $\epsilon_i\sim N(0,\sigma^2)$. The sample mean is $\bar{X}=\mu+(1/n)\sum\epsilon_i$, so $\E[\bar{X}]=\mu+0=\mu$, and $\Var(\bar{X})=(1/n^2)\cdot n\sigma^2=\sigma^2/n$.
		\item The standard error is $\text{SE}(\bar{X})=\sqrt{\Var(\bar{X})}=\sigma/\sqrt{n}$. To reduce the error by half, $n$ must be quadrupled (the dependence is through $\sqrt{n}$).
		\item For $\sigma/\sqrt{n}<0.01$ with $\sigma=0.1$: $\sqrt{n}>10$, so $n>100$. With $n=100$ measurements, the standard error is exactly $0.01$ (1\% of the true value).
	\end{enumerate}
\end{solproblema}

% Analyze
\begin{solproblema}[prob:de8d740]
	By the law of total expectation, $\E[X]=\E[\E[X\mid Y]]$. Expand $\Var(X)$:
	\begin{align*}
		\Var(X) &= \E[(X-\E[X])^2] = \E\big[\E[(X-\E[X\mid Y]+\E[X\mid Y]-\E[X])^2 \mid Y]\big] \\
		&= \E[\E[(X-\E[X\mid Y])^2\mid Y]] + \E[\E[(\E[X\mid Y]-\E[X])^2\mid Y]] \\
		&\quad + 2\E\big[\E[(X-\E[X\mid Y])(\E[X\mid Y]-\E[X])\mid Y]\big].
	\end{align*}
	The third term is zero by the tower property: $\E[X-\E[X\mid Y]\mid Y]=\E[X\mid Y]-\E[X\mid Y]=0$. The first term is $\E[\Var(X\mid Y)]$ and the second is $\Var(\E[X\mid Y])$. Therefore:
	\begin{align*}
		\Var(X) = \E[\Var(X \mid Y)] + \Var(\E[X \mid Y]).
	\end{align*}
	In ANOVA, total variability decomposes into \emph{within-group variability} ($\E[\Var(X\mid Y)]$, analogous to SSE) and \emph{between-group variability} ($\Var(\E[X\mid Y])$, analogous to SSTR).
\end{solproblema}

% Evaluate
\begin{solproblema}[prob:9d4a41b]
	For the exponential distribution with $\lambda=1$: $\mu=1$, $\sigma=1$, $\E[(X-1)^3]=2$, $\E[(X-1)^4]=9$. Therefore:
	\begin{align*}
		\gamma_1 = 2/1^3 = 2 \quad \text{(positive skewness)}, \qquad \gamma_2 = 9/1^4-3 = 6 \quad \text{(excess kurtosis)}.
	\end{align*}
	Compared with the standard Normal distribution ($\gamma_1=0$, $\gamma_2=0$), the exponential has marked positive skewness: its right tail is much longer than its left side (which does not even exist, since the support is $[0,\infty)$), reflecting that values far above the mean, while rare, are substantially more likely than in a symmetric distribution. The positive excess kurtosis ($\gamma_2=6 \gg 0$) indicates a \emph{leptokurtic} distribution: more peaked near the origin and with heavier tails than the Normal, meaning greater probability of both very small values (near zero) and extremely large values compared with what a normal approximation would predict.
\end{solproblema}

% Create
\begin{solproblema}[prob:5c186f2]
	\textbf{Example:} let a portfolio have expected returns $\E[R_1]=0.08$ (8\%, stocks), $\E[R_2]=0.05$ (5\%, bonds), $\E[R_3]=0.12$ (12\%, high-risk fund), with weights $w_1=0.40$, $w_2=0.35$, $w_3=0.25$.

	By linearity of expectation:
	\begin{align*}
		\E[w_1R_1+w_2R_2+w_3R_3] = w_1\E[R_1]+w_2\E[R_2]+w_3\E[R_3] = 0.40(0.08)+0.35(0.05)+0.25(0.12).
	\end{align*}
	\begin{align*}
		= 0.032+0.0175+0.030 = 0.0795 \text{ (7.95\%).}
	\end{align*}
	Linearity of expectation does not require independence among $R_1, R_2, R_3$: the general proof $\E[aX+bY]=a\E[X]+b\E[Y]$ relies only on the linearity of the integral (or sum) that defines expectation, without using at any point the factorization of a joint density that independence would require. Therefore, the portfolio's expected return is valid and computable this way even if, in practice, the three assets are correlated (as often happens in real financial markets).
\end{solproblema}
